\section{Pollard rho}
\label{sec:rho}

Pollard rho is the reference. This section records the walk this
library actually runs, the constants it measures, and the GPU and
fleet machinery that parallelise it without changing $S$.

\subsection{The walk}
\label{sec:rho-walk}

The solver is Teske's $r$-adding walk~\cite{teske2001}. Multipliers
$M_i=\alpha_i G+\beta_i H$ are drawn once per job; the state is
$(Y,a,b)$ with $Y=aG+bH$; a step is $Y\leftarrow Y+M_{\mathrm{idx}(Y)}$.
Distinguished points --- hashes with \texttt{dp\_bits} trailing zero
bits --- are stored with their $(a,b)$ in a shared table (van
Oorschot--Wiener~\cite{vanoorschot1999}). A second arrival with a
different $(a,b)$ gives $(b-b')x=a'-a\pmod n$. For composite $n$ the
congruence is solved by gcd and every candidate is
verified.\art{src/rho.c}\art{docs/ALGORITHMS.md, ``Pollard rho''}

$r$, the number of walks per thread $W$, and \texttt{dp\_bits} are
chosen automatically from $n$ and the thread count so that neither the
multiplier table nor the walk start-ups (each about $3\log_2 n$
operations) exceed a small fraction of $\sqrt{n}$, and so that the
expected number of distinguished points is about 32 per walk, bounded
above at $2^{24}$ table entries. Walks that run $24\cdot
2^{\texttt{dp\_bits}}$ steps without a distinguished point are
restarted.

Threads run batches of $W$ walks; each batch step is one
\texttt{ca\_group\_batch\_op}, Montgomery's simultaneous
inversion~\cite{montgomery1987}. On one core an affine curve addition
costs $232.2$~ns with a private inversion and $21.8$~ns in batches of
$256$ ($45.88$~Mops/s), a $10.6\times$
saving.\art{docs/BENCHMARKS.md, ``Raw group operation throughput''} A
$61$-bit $(\ZZ/p\ZZ)^*$ Montgomery multiply is $5.9$~ns
($169.4$~Mops/s) on the same host.

\subsection{Negation map and fruitless cycles}

On elliptic curves every point is replaced by the canonical member of
$\{Y,-Y\}$ (the one with the smaller Montgomery $y$), negating $(a,b)$
when needed. The search space shrinks by $2$ and the expected cost by
$\sqrt{2}$. The walk on classes is well defined because $\mathrm{idx}$
is evaluated on the canonical form.

Fruitless cycles are handled as in Bernstein--Lange--Schwabe~\cite{bls2011}:
\begin{itemize}[leftmargin=1.2em]
\item \emph{2-cycles.} The look-ahead rule ``if $\mathrm{idx}(Y+M_i)=i$
  use $M_{i+1}$ instead'' (iterated) keeps the walk a deterministic
  function of $Y$ and reduces the 2-cycle rate from $1/(2r)$ to about
  $1/(2r)^2$.
\item \emph{Longer cycles.} A sliding window compares each new point
  with a point saved every 64 steps. On a hit the cycle is traversed
  once, its minimum (by hash) found, and the walk escapes
  deterministically to $2\cdot\min$. Because the escape depends only
  on the cycle, two walks that merged inside a cycle stay merged.
\end{itemize}
With the negation map, $r$ is allowed to grow to $1024$; without it
the default is $32$.

\subsection{Measured $S$}
\label{sec:rho-S}

Table~\ref{tab:rho-s} is the whole-group table of
\texttt{docs/BENCHMARKS.md}: prime-order $(\ZZ/p\ZZ)^*$ subgroups of
safe primes and prime-order curves, five instances per row, $S$
counting every thread's operations. With five repetitions the rho and
kangaroo rows carry about $\pm 25\%$ statistical noise (the standard
deviation of one run is about half its mean), so the honest summary is
that rho, kangaroo and grumpy giants all sit at $S\sim 1.2$--$1.8$,
BSGS at $1.5$, as theory says.\art{docs/BENCHMARKS.md, ``Whole-group discrete logarithm''}

\begin{table}[t]
\centering
\small
\caption{Whole-group discrete logarithm, $S=\mathrm{ops}/\sqrt{n}$,
  5 instances per row, 4 threads for \texttt{rho-T}. From
  \texttt{ca\_bench generic --threads 4}.}
\label{tab:rho-s}
\begin{tabular}{@{}lrrrrr@{}}
\toprule
grp & bits & BSGS & rho-1 & kangaroo & grumpy \\
\midrule
$\ZZ_p^*$ & 24 & 1.496 & 2.066 & 1.751 & 1.242 \\
$E(\FF_p)$ & 24 & 1.496 & 1.603 & 1.886 & 1.211 \\
$\ZZ_p^*$ & 32 & 1.575 & 1.256 & 1.412 & 1.365 \\
$E(\FF_p)$ & 32 & 1.575 & 0.791 & 2.174 & 1.514 \\
$\ZZ_p^*$ & 36 & 1.441 & 1.563 & 1.691 & 1.351 \\
$E(\FF_p)$ & 36 & 1.441 & 1.283 & 1.768 & 1.413 \\
$\ZZ_p^*$ & 40 & 1.573 & 0.887 & 1.795 & 1.393 \\
$E(\FF_p)$ & 40 & 1.573 & 1.713 & 1.220 & 1.251 \\
\bottomrule
\end{tabular}
\end{table}

One-thread rho on $32$-bit $\ZZ_p^*$ is $S=1.256$; on curves with the
negation map the 24--40-bit span is $0.791$--$1.713$. The wall-clock
column of the source table (not reproduced) shows the structural
facts: table-based methods slow down per operation as the table grows,
rho stays at $40$--$100$~Mops/s in $\ZZ_p^*$ and $10$--$20$~Mops/s on
curves, and kangaroo is the fastest per operation because it stores
almost nothing.

Interval logarithms of width $2^w$ inside a $\sim 2^{60}$ ($\ZZ_p^*$)
or $\sim 2^{56}$ (curve) group sit at $S=\mathrm{ops}/\sqrt{w}$ of
$1.5$--$2.8$. A separate 400-instance run of grumpy giants gives
$S=1.78$ for intervals and $1.18$ for whole groups; the five-instance
rows scatter around those
values.\art{docs/ALGORITHMS.md, ``Two grumpy giants and a baby''}

\subsection{Fitted exponents}

$S$ assumes the exponent is $1/2$. The library also measures it.
\texttt{ca\_bench complexity --bits 20,24,28,32,36 --reps 5} fits
$\mathrm{cost}\sim C\cdot N^{\alpha}$ in log-log space and reports
$\alpha$ next to the theoretical exponent
(Table~\ref{tab:exponents}).\art{docs/BENCHMARKS.md, ``Measured time complexity''}

\begin{table}[t]
\centering
\small
\caption{Fitted exponents, five sizes, five repetitions. The last
  column is $\mathrm{mean}(\mathrm{cost}/N^{\mathrm{theory}})$, the
  familiar $S$, meaningful because the fitted $\alpha$ confirms the
  exponent.}
\label{tab:exponents}
\begin{tabular}{@{}llrrr@{}}
\toprule
algorithm & theory & $\alpha$ & $R^2$ & const \\
\midrule
BSGS ($\ZZ_p^*$) & $N^{0.500}$ & 0.487 & 0.9993 & 1.508 \\
rho ($\ZZ_p^*$) & $N^{0.500}$ & 0.477 & 0.9903 & 1.781 \\
kangaroo ($\ZZ_p^*$) & $N^{0.500}$ & 0.526 & 0.9910 & 1.843 \\
grumpy ($\ZZ_p^*$) & $N^{0.500}$ & 0.502 & 0.9966 & 1.264 \\
precomp build ($\ZZ_p^*$) & $N^{0.667}$ & 0.647 & 0.9954 & 1.192 \\
precomp online ($\ZZ_p^*$) & $N^{0.333}$ & 0.347 & 0.8768 & 2.781 \\
BSGS (EC) & $N^{0.500}$ & 0.487 & 0.9993 & 1.508 \\
rho (EC) & $N^{0.500}$ & 0.428 & 0.9845 & 1.572 \\
kangaroo (EC) & $N^{0.500}$ & 0.494 & 0.9995 & 1.622 \\
grumpy (EC) & $N^{0.500}$ & 0.480 & 0.9956 & 1.128 \\
\bottomrule
\end{tabular}
\end{table}

The square-root methods land at $\alpha\sim 0.5$ and the
precomputation phases at $0.667$ and $0.333$. The EC-rho fit
$\alpha=0.428$ is a finite-size artefact (negation-map savings and
walk-start constants on 20--36-bit curves), not a claim that rho
beats $N^{1/2}$. Index calculus is absent from the table on purpose:
it is $L_p[1/2]$, so no single power-law describes it
(Section~\ref{sec:ic-mult}).

\subsection{GPU kernel, emulator-verified}
\label{sec:rho-gpu}

The same walk runs as a CUDA kernel: tens of thousands of walks, one
thread owning eight of them that share a single Fermat inversion
(\texttt{CA\_GPU\_W}${}=8$). A private inversion would be about $99$
multiplications per curve step; the batched step is about $18$, a
$5.5\times$ arithmetic saving.\art{docs/GPU.md}

The device code is simultaneously valid CUDA~C++ and valid C11. A host
emulator executes the identical kernel body one thread at a time, so
the algorithm is testable without a GPU. Table~\ref{tab:gpu-emu} is
that emulator: operation counts are portable, seconds are not a GPU
measurement, and no row is a device-throughput
result.\art{docs/BENCHMARKS.md, ``GPU rho kernel''}

\begin{table}[t]
\centering
\small
\caption{Host-emulator GPU rho versus CPU rho. $S$ does not grow with
  the walk count: van Oorschot--Wiener parallelisation is linear in
  the number of walkers. 3 instances per row.}
\label{tab:gpu-emu}
\begin{tabular}{@{}lrrrr@{}}
\toprule
grp & bits & CPU $S$ & GPU-emu $S$ & walks \\
\midrule
$\ZZ_p^*$ & 24 & 1.849 & 1.895 & 8 \\
$E(\FF_p)$ & 24 & 1.707 & 1.844 & 8 \\
$\ZZ_p^*$ & 28 & 2.578 & 1.990 & 8 \\
$E(\FF_p)$ & 28 & 2.033 & 1.504 & 16 \\
$\ZZ_p^*$ & 32 & 2.045 & 1.262 & 40 \\
$E(\FF_p)$ & 32 & 0.873 & 1.857 & 56 \\
\bottomrule
\end{tabular}
\end{table}

Both solvers sit at $S\sim 1.3$--$2.6$. A regression --- an $S$ that
grows with the walk count --- would mean the distinguished-point
density or the restart policy is wrong, which is what this row
watches. Kernel resources from \texttt{ptxas} 12.9, no GPU required:
72 registers and $736$~B local memory per thread on sm\_70, 64
registers and the same local memory on sm\_80/89/90, zero spills on
all four.\art{docs/GPU.md, resource table}

The host never reads walk state. Distinguished points are written
through \texttt{atomicAdd}; the host verifies every candidate with a
scalar multiplication before returning it.

\subsection{Fleet protocol}
\label{sec:rho-fleet}

\texttt{ca\_coord} divides one instance across machines that cannot
reach each other. Agents in private subnets dial a coordinator URL;
the hub pushes distinguished points and the solution back down the
socket the agent opened. What is shared is van Oorschot--Wiener's
observation that parallel rho divides by start point, not by search
space: every agent walks the same function (the branch table is
derived from the job id) and reports only distinguished points. What
is partitioned is the walker index space. Every fact on the wire is
verifiable from the job alone: a check-in carries $({\mathrm{walker}},
{\mathrm{steps}}, Y, a, b)$ with $aG+bH=Y$, which costs the receiver
two scalar multiplications to check against the $2^{\texttt{dp\_bits}}$
steps it stands for. The hub is a rendezvous, not an
authority.\art{include/cryptanalysis/ca_coord.h}

The ECC2K-130 campaign of Section~\ref{sec:campaign} is the same
protocol specialised to a Koblitz iteration: 16-bit \texttt{--run-id}
spaces, 32-byte distinguished-point records, checkpointed iteration
totals rather than point counts for rate estimates, and a public
dashboard that publishes aggregates only.
